\setchoices{label=(\Alph*)}

\begin{question}[index=1]
函数 $y=y(x)$ 由方程 $|x|y+\eu^y=\eu$ 确定，则下列关于点 $x=0$ 的说法正确的是
\begin{choices}[columns=2]
\item 不可导点，$y(0)$ 是极大值
\item 不可导点，$y(0)$ 不是极值
\item 可导点，$y(0)$ 不是极值
\item 可导点，$y(0)$ 是极大值
\end{choices}
\end{question}
\QuestionBreak

\begin{question}[index=2]
记 $L$ 为曲线 $\begin{cases}x^2+y^2+z^2=1,\\2x+2y+3z=0,\end{cases}$ 则曲线积分
$\displaystyle I=\int_L(3x^2+y^2+2z^2+x+2z)\dd s=$
\begin{choices}[columns=4]
\item $\pi$
\item $2\pi$
\item $3\pi$
\item $4\pi$
\end{choices}
\end{question}
\QuestionBreak

\begin{question}[index=3]
级数 $\displaystyle\sum_{n=2}^{\infty}\frac{x^{2n}}{n\ln^2 n}$ 的收敛域为
\begin{choices}[columns=4]
\item $(-1,1)$
\item $[-1,1]$
\item $(-\sqrt{2},\sqrt{2})$
\item $[-\sqrt{2},\sqrt{2}]$
\end{choices}
\end{question}
\QuestionBreak

\begin{question}[index=4]
函数 $\displaystyle y=\int_0^x\frac{\left(t-\frac{\pi}{2}\right)^{\frac43}}{1-\sin^3t}\dd t$ 的定义域为
\begin{choices}[columns=4]
\item $\left(-\dfrac{3\pi}{2},\dfrac{5\pi}{2}\right)$
\item $\left(-\dfrac{\pi}{2},\dfrac{\pi}{2}\right)$
\item $\left(-\dfrac{3\pi}{2},\dfrac{\pi}{2}\right)$
\item $\left(-\dfrac{3\pi}{2},\dfrac{3\pi}{2}\right)$
\end{choices}
\end{question}
\QuestionBreak

\begin{question}[index=5]
设 $3$ 阶行列式 $|\vect A|$ 中各列元素之和为 $2$，第一行元素的代数余子式之和 $A_{11}+A_{12}+A_{13}=2$，
$\vect A^*,\vect A^{-1}$ 分别为 $\vect A$ 的伴随矩阵和逆矩阵，则 $|\vect A^*-2\vect A^{-1}|=$
\begin{choices}[columns=4]
\item $0$
\item $1$
\item $2$
\item $4$
\end{choices}
\end{question}
\QuestionBreak

\begin{question}[index=6]
设 $\vect\alpha_1,\vect\alpha_2,\ldots,\vect\alpha_s$ 为线性无关的 $n$ 维列向量组，$\vect A$ 为 $n$ 阶矩阵，
向量组 $\vect A\vect\alpha_1,\vect A\vect\alpha_2,\ldots,\vect A\vect\alpha_s$ \allowbreak\mbox{线性相关是 $r(\vect A)<n$ 的}
\begin{choices}[columns=2]
\item 充分必要条件
\item 充分非必要条件
\item 必要非充分条件
\item 非充分非必要条件
\end{choices}
\end{question}
\QuestionBreak

\begin{question}[index=7]
二次型 $f(x_1,x_2,x_3)=(x_1+x_2-x_3)(x_1+2x_2+3x_3)$ 的规范型为
\begin{choices}[columns=4]
\item $y_1^2+y_2^2$
\item $y_1^2-y_2^2$
\item $y_1^2-y_2^2+y_3^2$
\item $y_1^2+y_2^2+y_3^2$
\end{choices}
\end{question}
\QuestionBreak

\begin{question}[index=8]
设 $P(A)=0.5$，$P(B)=0.8$，则 $P(\overline A\,\overline B)P(B\mid\overline A)$ 的最大值为
\begin{choices}[columns=4]
\item $0.125$
\item $0.12$
\item $0.08$
\item $0.15$
\end{choices}
\end{question}
\QuestionBreak

\begin{question}[index=9]
设随机变量 $X,Y$ 相互独立，且 $X\dist N(1,1)$，$Y\dist N(2,1)$，且 $D[(2X-Y)^2]=$
\begin{choices}[columns=4]
\item $50$
\item $10$
\item $18$
\item $9$
\end{choices}
\end{question}
\QuestionBreak

\begin{question}[index=10]
设随机变量 $(X,Y)$ 的联合密度函数为 $f(x,y)=\begin{cases}\eu^{-x},&0<y<x,\\0,&\text{其他},\end{cases}$ 则 $P\{X>2\mid Y=1\}=$
\begin{choices}[columns=4]
\item $\eu^{-2}$
\item $\eu^{-1}$
\item $1-\eu^{-1}$
\item $1-\eu^{-2}$
\end{choices}
\end{question}
